\scanpage{32}
\mspage{\textcircled{26}}

Conditions sur $(M,W)$ avec \deleted{\unsure{des}}

\deleted{\unsure{(1)}} a) $M$ stable par $\varprojlim$ finies et $\varinjlim$ (petites).

b) $W$ \deleted{\unsure{il en est}} stable par $\varinjlim$ filtrantes

c) $M$ accessible.

\diagram{\scriptsize\begin{tikzcd}[ampersand replacement=\&,row sep=small,/tikz/scale=.82,/tikz/transform shape]
 {\mathrm I\left\{\begin{aligned}
 &\text{\textcircled{a} existence d'un syst. de Quillen faible }(W,C,TC,F,TF)\ \text{(clos si on y tient)}\\
 &\Updownarrow\\
 &\text{\textcircled{b} existence d'un syst. de Quillen faible \unsure{ex.ct|ex.t|exact} à gauche }(TC=C\cap W)\\
 &\Updownarrow\\
 &\text{\textcircled{c} }\underbrace{(\mathrm{Cof}_W)_*\subset W,\ \deleted{\unsure{\mathrm{Cof}}}(W^{\mathrm{univ}})_*\subset\mathrm{Fib}_W}_{(*)}
 \ \Longleftrightarrow\ (\mathrm{Cof}_W)_*\subset W_{\mathrm{univ}},\ (W^{\mathrm{univ}})_*\subset\mathrm{Fib}_W
 \end{aligned}\right.}\\
 {\mathrm{II}\left\{\begin{aligned}
 &\text{\textcircled{a} existence d'un syst. de Quillen faible \unsure{ex.t|ex.ct|exact} à droite }(TF=F\cap W)\\
 &\Updownarrow\\
 &\text{\textcircled{\unsure{d|b}} On a }(*),\ \text{et de plus }(W^{\mathrm{univ}})_*\cap W_{\mathrm{univ}}\ \text{clos à droite}\\
 &\Updownarrow\\
 &\deleted{\unsure{\mathrm{III}}}\ \text{\textcircled{b} existence d'un syst. de Quillen clos (closed model category) :}\\
 &\qquad TC=C\cap W,\ TF=F\cap W,\ F=(TC)_*,\ TF=C_*,\ TC=F^*,\ C=(TF)^*\\
 &\tikz[baseline=-.5ex]{\draw[<->,double] (-.19,.13)--(.19,-.13);}\\
 &\text{\textcircled{c} Posant }TC=W^{\mathrm{univ}},\ F=(TC)_*,\ TF=F\cap W_{\mathrm{univ}},\ C=(TF)_*,\\
 &\qquad\text{on trouve avec }W\text{ un syst. de Quillen clos}
 \end{aligned}\right.}\arrow[u,Rightarrow]\\
 {\mathrm{III}\left\{\begin{aligned}
 &C=\mathrm{Cof}_W\ \text{et }TC=W^{\mathrm{univ}}\ \text{font partie \inserted{(avec $W$)} d'un syst. de Quillen clos}\\
 &\qquad(\text{donc }F=(W^{\mathrm{univ}})_*,\ TF=(\mathrm{Cof}_W)_*,\ W=TF\circ TC\unsure{)}\\
 &\Updownarrow\\
 &(*),\ \text{et de plus }(W^{\mathrm{univ}})_*\cap W_{\mathrm{univ}}=(\mathrm{Cof}_W)_*
 \end{aligned}\right.}\arrow[u,Rightarrow]\\
 {\mathrm{IV}\left\{\begin{aligned}
 &\text{On a }(*)\text{ et }\mathrm{Cof}_W=\widetilde{(W^{\deleted{\unsure{\mathrm{univ}}}}_{\mathrm{univ}})^*\cup W_{\deleted{\unsure{\mathrm{univ}}}}^{\mathrm{univ}}}\\
 &\Updownarrow\\
 &\text{On a }(*)\text{ et }(W^{\mathrm{univ}})_*\cap\overline{W_{\mathrm{univ}}}=(\mathrm{Cof}_W)_*
 \end{aligned}\right.}\arrow[u,Rightarrow]
\end{tikzcd}}

Cas $M=\mathrm{Cat}$, $W=W_\infty$. On sait que \unsure{II|III} est satisfait, par
Thomason. Donc on a \deleted{\unsure{III} \illegible{1}}
\[
 (\mathrm{Cof}_W)_*\subset W_{\mathrm{univ}},\quad(W^{\mathrm{univ}})_*\subset\mathrm{Fib}_W,\quad
 (W^{\mathrm{univ}})_*\cap W_{\mathrm{univ}}\ \text{clos à \unsure{droite}}
\]
J'ai pu prouver directement les deux premières inclusions
(i.e. $\mathrm I$), mais pas que $(W^{\mathrm{univ}})_*\cap W_{\mathrm{univ}}$ soit clos
à droite. (Je devrais au moins essayer !)

\deleted{\unsure{Je} \unsure{signale que} \illegible{2} \unsure{du tout} \unsure{on a}} \inserted{On n'a \uuline{pas} (IV), i.e. on a}
\[
 (W^{\mathrm{univ}})_*\cap\overline{W_{\mathrm{univ}}}\ne(\mathrm{Cof}_W)_*
 \deleted{\unsure{(W^{\mathrm{univ}})_*\cap\overline{W_{\mathrm{univ}}}}}
\]
car
\[
 (W^{\mathrm{univ}})_*\cap\overline{W_{\mathrm{univ}}}\not\subset W_{\mathrm{univ}}
\]
(contient les $X\to e$ avec $\unsure{\alpha|X}$ discrète)
donc on a
\[
 (W^{\mathrm{univ}})_*\cap\overline{W_{\mathrm{univ}}}\supsetneq
 \underbrace{(W^{\mathrm{univ}})_*\cap W_{\mathrm{univ}}}_{\substack{\text{\unsure{n'est pas l'orth.} de}\\W^{\mathrm{univ}}\cup(W_{\mathrm{univ}})^*}}
\]
\margin{Il suffirait de savoir que $\overline{W_{\mathrm{univ}}}\cap\mathrm{Fib}_W=W_{\mathrm{univ}}$ ce qui impliquerait que III et IV sont équivalents}
\inserted{Mais c'est faux, p.ex. si $X$ cat. discr. non vide, $X\to e$ est dans $\overline{W_{\mathrm{univ}}}\cap\mathrm{Fib}_W$}
% archive sheet number: 31
